\section{Évolution vers la robotique collaborative}

Après la publication de l’article initial, les travaux ont évolué vers une analyse approfondie de la robotique collaborative.

Les systèmes HRC impliquent :
\begin{itemize}
    \item la co-présence humain–robot,
    \item la coordination de tâches,
    \item la gestion des distances de sécurité,
    \item l’interprétation d’intentions humaines.
\end{itemize}

Les revues structurantes du domaine \cite{hrc_synthesis} mettent en évidence plusieurs niveaux de collaboration : coexistence, coopération et collaboration complète.

Les travaux récents en interaction cognitive et sociale \cite{access_2024} soulignent l’importance de modèles adaptatifs intégrant perception multimodale, raisonnement contextuel et apprentissage dynamique.

Ces analyses confirment que la collaboration homme–robot dépasse la simple planification de trajectoire et nécessite des mécanismes d’adaptation comportementale intégrés.
